\documentclass{article}
\usepackage{amsmath}
\newcommand{\lagL}{\mathcal{L}}
\newcommand{\ordO}{\mathcal{O}}
\newcommand{\actS}{\mathcal{S}}
\newcommand{\tS}{\tilde{S}}
\newcommand{\vx}{\vec{x}}
\newcommand{\vk}{\vec{k}}
\newcommand{\vq}{\vec{q}}
\newcommand{\vQ}{\vec{Q}}
\newcommand{\vr}{\vec{r}}
\newcommand{\vA}{\vec{A}}
\newcommand{\vS}{\vec{S}}
\newcommand{\vL}{\vec{L}}
\newcommand{\vN}{\vec{N}}
\newcommand{\un}{\hat{n}}
\newcommand{\ue}{\hat{e}}
\begin{document}

\title{The effect of bond impurity in triangular
lattice heisenberg anti-ferromagnet}
\maketitle

\begin{abstract}
    In this note we explore the fate of the 120 
    degree long range order in the presence of
    randomly distributed bond impurities.
\end{abstract}

\section{Introduction}
In the context of the classical spiral ground 
state of the triangular lattice heisenberg anti
-ferromagnet, the presence of a single bond 
impurity locally releases the frustration and 
helps the associated spins to anti-align. In 
the reference frame of an individual spin this
amounts to a perturbation in the transverse
direction of the magnetic order. If we assume 
a rotated co-ordinate system for each spin such
that in the absence of any fluctuation, 
$\langle \tS^z\rangle=S$, the effect of
a single impurity could be understood as a 
perturbation to the hamilonian,
\begin{equation}
    \begin{aligned}
        \delta H_i =h_i\sum_{j\epsilon\{i\}}
        \alpha_j\tS^x_j
        =h_iS^{\text{imp}}_i
    \end{aligned}
\end{equation}
where the notation is as follows; the bond 
impurity has been introduced between the coupling
of the $i$'th spin and one of its neighbour 
and $j$ is either 0 or 1 depending on whether
it designates the $i$'th spin or its neighbour. 
The phase factor $\alpha_j=\pm(-1)^{j}$ ensures 
the correct symmetry of the perturbation. Because
of the spiral pattern of the underlying spin
pattern the neighbouring spins have to rotate 
in the opposite directions in-order to anti-align.
The overall $\pm$ of the phase factor is fixed
with the breaking of the chiral symmetry of the
$120^\circ$ order. However, this does not have
any effect on the observables that we are interested
in.

\section{The response function}
Within the linear response approximation the transverse
spin fluctuation of an arbitrary spin $j$ away from the 
impurity at $i$ would be given by,
\begin{equation}
    \begin{aligned}
        \langle\tS_j^x\rangle &= h_{i}
        \langle\tS^x_jS_i^{\text{imp}}\rangle
        =h_i\langle\tS^x_j\sum_{j'\epsilon\{i\}}
        \alpha_{j'}\tS^x_{j'}\rangle \\
        &=h_i\sum_{\vk,\vk'}
        \langle\tS^x_{\vk}
        \left(\sum_{j'\epsilon{i}}\alpha_{j'}
        e^{i\vk'\cdot\vr_{j'}}\right)
        \tS^x_{\vk'}
        \rangle 
        e^{i\vk\cdot\vr_j} \\
        &=h_i\sum_{\vk}\alpha_{\vk}
        \langle\tS^x_{\vk}\tS^x_{-\vk}\rangle
        e^{i\vk\cdot(\vr_j-\vr_i)} \\
        &=h_iG^\perp(\vr_j-\vr_i)
    \end{aligned}
\end{equation}
Where we use the definition, 
$\alpha_{\vk}=(-1)^0+(-1)^1e^{-i\vk\cdot\vec{\delta}}$,
with $\vec{\delta}$ being the bond directional vector. 
The choice of any of the six directional vectors around 
a particular spin should not affect the overall result.
In the following we choose the direction 
$\vec{\delta}=(1,0)$ for our calculation. Following
Wollny, the momentum space response function at $T=0$
energy can now simply be written as,
\begin{equation}
    \begin{aligned}
        &G^\perp({\vk}) = \alpha_{\vk}
        \langle\tS^x_{\vk}\tS^x_{-\vk}\rangle
        =S\frac{\alpha_{\vk}}{\gamma_{\vk}-1} \\
        &\text{ with,} \\
        &\alpha_{\vk} = 1-e^{-ik_x} \\
        &\gamma_{\vk} = 
        \frac{1}{3}\left(\cos k_x
        +2\cos(k_x/2)\cos(\sqrt{3}k_y/2)\right)
    \end{aligned}
\end{equation}
The principle contributions for the response function 
in the zero energy limit will come from the goldstone
moders around the ordering wave vectors for triangular
lattice, $\vQ=(0,0) \text{ and } \pm(4\pi/3,0)$. 
We expand around these points, $\vk=\vq+\vQ$ with a
polar parametrization $\vq=(q\cos\phi,q\sin\phi)$ and
find,
\end{document}
